\subsection{Countable trace completions: an upper bound and a duality obstruction}
The finite-spectrum hypothesis has a useful extension, but its duality
requirement becomes restrictive. Call a visible weighted spectrum
$\mathcal A=\{(\alpha,c_\alpha):\alpha\in A\}$ \emph{admissible} if
$A\subset\C^*$ is bounded, its roots are distinct, every
$\{\alpha:|\alpha|\ge\epsilon\}$ is finite, all $c_\alpha\ne0$, and
for some nonnegative integer $N$,
\[
 S_N=\sum_{\alpha\in A}|c_\alpha|\,|\alpha|^N<\infty.
\]
Equal roots are grouped and zero coefficients removed before imposing
these conditions. The empty spectrum represents zero.

\begin{finding}[\fid{F32}: admissible countable spectra and visible reciprocal duality]\label{f:countable}
For a nonempty admissible spectrum, the maximum radius $R$ is attained
and $u_n=\sum c_\alpha\alpha^n$ converges absolutely for $n\ge N$.
It satisfies
\[
 \limsup_{n\to\infty}|u_n|^{1/n}=R.
\]
An eventual sequence has at most one admissible visible expansion.
Finite sums, tensor products, and integer power substitutions preserve
admissibility after grouping, although power substitution may erase the
entire largest-radius part.

The signed even-path criterion of \fid{F24} remains true when its
input visible spectrum and each positive path-output spectrum are
admissible countable spectra. The identities and bounds need hold only
for all sufficiently large extension degrees, with thresholds depending
on the path length. The input coefficients must still be real, and the
finite-extension and all-embedding versions retain the same explicit
visibility and transformed-hypothesis requirements.

If $A$ is admissible and $\kappa\ne0$, and every $\kappa/\alpha$ is
visible in a bounded dual support, then $A$ is finite. In particular,
an infinite completion of this kind cannot have bounded visible
reciprocal duality. If the nonzero grouped coefficients are integers,
absolute convergence of one nonnegative moment already implies the
boundedness and local-finiteness conditions above.
\end{finding}
\begin{proof}
Choose a root $\alpha_0$. The finite set of roots of modulus at least
$|\alpha_0|$ contains a largest one; it is a global maximum $R>0$.
For $n\ge N$, absolute summation gives $|u_n|\le S_NR^{n-N}$.
The generating function of the shifted sequence is
\[
 H_N(T)=\sum_{m\ge0}u_{N+m}T^m
       =\sum_{\alpha\in A}\frac{c_\alpha\alpha^N}{1-\alpha T}
       \quad (|T|<1/R).
\]
The equality follows from the double absolute bound
$S_N/(1-R|T|)$. Its right side converges locally normally away from
$\{1/\alpha\}$. Indeed, on a compact set with $|T|\le L$, roots
of modulus at most $1/(2\max(1,L))$ have denominator modulus at
least $1/2$, giving a normally convergent tail by $2S_N$; only
finitely many other roots remain. The pole set is locally finite.
At $T=1/\alpha$ the residue is
$-c_\alpha\alpha^{N-1}\ne0$; all other terms are holomorphic in a
neighborhood. Thus the Taylor series has radius exactly $1/R$.
Cauchy--Hadamard, and the harmless fixed shift by $N$, give the
asserted exponential rate. To prove uniqueness, subtract two proposed
expansions at a common sufficiently large moment index. Their grouped
difference is admissible. If nonempty its exponential rate is positive,
contradicting an eventually zero sequence.

For a tensor power of order $k$, a product of modulus at least
$\epsilon$ has every factor of modulus at least
$\epsilon/R^{k-1}$. Only finitely many factor tuples can therefore
produce any nonzero product, or any product outside a prescribed disk.
After grouping to coefficients $d_{k,\rho}$,
\[
 \sum_\rho |d_{k,\rho}|\,|\rho|^N\le S_N^k.
\]
This proves tensor closure, including all multiplicative relations
among roots. Finite sums use a common moment index and the triangle
inequality. For the substitution $n=er$, group the at most $e$ roots
with a given $e$th power and choose $N'$ with $eN'\ge N$; boundedness
and the moment estimate prove admissibility. Cancellation on the top
circle is possible, as $q^n-(-q)^n$ already shows.

For completeness, consider the countable signed path argument directly.
Write $w=a+d$ and let $A_R$ be the finite top circle. For every allowed
even $k$, the finite group-algebra element
$F=\sum_{\alpha\in A_R}c_\alpha[\alpha]$ has
$F^k\ne0$: otherwise evaluation $[\alpha]\mapsto\alpha^n$ for
all positive integers $n$, followed by Vandermonde invertibility,
would force every coefficient to vanish. Its coefficient sum is
$(\sum c_\alpha)^k\ge0$, so at least one of its real grouped
coefficients is positive. The corresponding root has modulus $R^k$.
No product using a smaller-radius root reaches that modulus, and the
positive path-output spectrum cannot cancel its coefficient.

Write the assumed path output as
$q^nN_{k,n}=-\sum_\beta b_{k,\beta}\beta^n+E_{k,n}$, with
$b_{k,\beta}>0$ and
$|E_{k,n}|\le C_kn^{s_k}q^{n(kw/2+\eta_k)}$,
$\eta_k/k\to0$. The exact path identity gives
\[
 M_n^k+\sum_\beta b_{k,\beta}\beta^n=E_{k,n}-R_{k,n},\qquad
 |R_{k,n}|\le B^kq^{n(kw/2+d/2+1)}.
\]
The left side is admissible by closure and visibly contains the root
just found. Applying the proved spectral-radius formula at fixed $k$
gives
\[
 R^k\le q^{kw/2+\max(\eta_k,d/2+1)}.
\]
Taking $k$th roots and then unbounded even $k$ proves
$R\le q^{w/2}$. Constants depending on $k$ and polynomial factors
in $n$ disappeared in the preceding fixed-$k$ limit. If only
$n=e_kr$ are used, the top circle must survive grouping for an unbounded
allowed set of $k$; then the same estimate has the factor $e_k$ on
both sides and yields the same conclusion. For each field embedding,
admissibility, real coefficients, estimates and visibility must be
verified for the transformed system separately.

For the reciprocal assertion let the dual support have radius bound
$R_D<\infty$. It is positive when $A\ne\varnothing$, and
$|\alpha|\ge |\kappa|/R_D$ for every $\alpha\in A$.
Local finiteness away from zero makes $A$ finite. No convergence
assumption on the dual support is needed for this deduction.
Consequently the two independent bounds used in \fid{F17}, together
with visible reciprocity $q^d/A$, still give purity, but now also
force the original visible support to be finite.

Finally suppose the coefficients are nonzero integers and the moment
sum converges. For $N=0$ every term is at least one, so the support is
finite. For $N\ge1$,
\[
 \#\{\alpha:|\alpha|\ge\epsilon\}\le S_N\epsilon^{-N},\qquad
 |\alpha|\le S_N^{1/N}.
\]
These prove the two support conditions. The requirement $N\ge0$
matters if boundedness was not separately given: roots $2^j$ with
coefficients one have a convergent moment at index $-1$.
\end{proof}

\paragraph{The natural completion of the actual normalized trace.}
The trace in \fid{F26}, after $v=q^{\pm n/2}$, is
\[
 h_n=24+\frac{16q^{-2n}}{(1+q^{-n})^4}
 =24+16\sum_{j\ge0}(-1)^j\binom{j+3}{3}
                   (q^{-(j+2)})^n.
\]
The binomial expansion is absolutely convergent for every $n\ge1$.
It gives an admissible integer Euler spectrum
$A=\{1,q^{-2},q^{-3},\ldots\}$, with
$S_1=24+16q^{-2}/(1-q^{-1})^4$. Its radius is one, consistent with
an upper weight-zero bound. This observation does not construct the
pointwise and path-output identities required by the theorem.
For every fixed $\kappa\ne0$, the reciprocal support
$\{\kappa,\kappa q^2,\kappa q^3,\ldots\}$ is unbounded.
With these nonzero integer coefficients its nonnegative moments do
not converge. Uniqueness excludes replacing the same eventual trace
by a different admissible finite or pure reciprocal expansion.
Thus this completion repairs the finite-spectrum requirement for an
upper-bound argument but cannot supply bounded visible reciprocal
duality for this actual word trace. Regularized traces, non-atomic
spectra, duals outside this convergence class, and other comparison
dictionaries are not excluded.
